% ============================================================
% WHITE PAPER TEXT — SIDM HALF-LIGHT RADIUS MASS FUNCTION
% ============================================================

% ── HIGH-LEVEL PARAGRAPH (for dark matter white paper body) ──────────────

\subsection*{Dark Matter Microphysics from the Enclosed Mass Function of Local Group Dwarfs}

The internal structure of dwarf spheroidal galaxies (dSphs) in the Local Group provides
one of the most sensitive probes of dark matter microphysics on sub-galactic scales.
Self-interacting dark matter (SIDM) predicts that elastic scattering between dark matter
particles thermalizes the inner regions of halos, driving them through a two-phase
gravothermal evolution: an initial core-forming phase in which energy is transported
outward and the central density decreases, followed by a core-collapsing phase in which
the inner halo becomes denser than its cold dark matter (CDM) counterpart.
The onset and duration of these phases depend sensitively on the effective
self-interaction cross section $\sigma_\mathrm{eff}/m$, which is itself a function of the
local dark matter velocity dispersion and hence of halo mass.
A powerful and observationally accessible statistic that captures this diversity is the
distribution of enclosed dark matter masses within stellar half-light radii,
$M(<r_{1/2})$, estimated via the Wolf et al.\ (2010) mass estimator from stellar
line-of-sight velocity dispersions.
We present a forward-model comparison of the $M(<r_{1/2})$ distribution for a
volume-complete sample of Local Group dSphs within 800~kpc against predictions from CDM,
SIDM, warm dark matter (WDM), and primordial power-spectrum bump models, and forecast
the discriminating power achievable with the combined Rubin Observatory photometric
survey and a next-generation wide-field spectrograph (SpecS5).


% ── DETAILED ANALYSIS PARAGRAPHS (for methods/supplementary section) ─────

\subsection*{Detailed Analysis Description}

\paragraph{Forward model.}
We construct a Monte Carlo population of $N = 2000$ dark matter halos by sampling
virial masses from the Tinker et al.\ (2008) analytic halo mass function (HMF) over the
range $10^8$--$10^{11}~M_\odot$ using inverse-CDF sampling via the \textsc{colossus}
package \citep{Diemer2018}.
All masses are defined with respect to the $200c$ overdensity threshold, consistent with
the Diemer \& Joyce (2019) mass--concentration relation used to assign NFW concentrations
with 0.15~dex log-normal scatter.
NFW structural parameters $(\rho_s, r_s)$ are extracted via the \textsc{colossus}
\texttt{NFWProfile} class, which handles the $h$-dependent unit conversions from
colossus internal units to physical units exactly.
Virial radii $R_\mathrm{vir}$ in the Bryan \& Norman (1998) overdensity definition are
also computed via \textsc{colossus} and used for the stellar size relation below.
Stellar half-light radii are assigned via a power-law size--halo relation
$r_{1/2} = 37 \times (R_\mathrm{vir}/10~\mathrm{kpc})^{1.07}$~pc with 0.63~dex
log-normal scatter, clipped to the observed dSph range of 20--3000~pc.
This parametrisation is based on the empirical scaling relation of
Jiang et al.\ (2019) and dominates the width of the predicted $M(<r_{1/2})$
distribution.

\paragraph{SIDM transformation.}
We apply the parametric SIDM model of Yang et al.\ (2023, hereafter Y23) to transform
each CDM halo into its SIDM counterpart.
For each halo, we compute the effective self-interaction cross section
$\sigma_\mathrm{eff}$ by numerically integrating the thermally-averaged Rutherford-like
differential cross section
$d\sigma/d\Omega = \sigma_0 w^4 / [2(w^2 + v^2\sin^2(\theta/2))^2]$
over the Maxwell--Boltzmann velocity distribution with characteristic velocity
$\nu_\mathrm{eff} = 1.1 V_\mathrm{max} / \sqrt{3}$,
where $V_\mathrm{max} = 1.648\,r_s\sqrt{G\rho_s}$ for an NFW halo.
We integrate up to $v_\mathrm{upper} = 10 \times \max(V_\mathrm{max}, w)$ to ensure
the cross section has turned over before truncation; using the smaller cutoff
$v_\mathrm{upper} = 3.06\,V_\mathrm{max}$ adopted by some earlier codes
underestimates $\sigma_\mathrm{eff}$ by factors of $\sim$2--10 when $V_\mathrm{max} \ll w$,
which is the relevant regime for dwarf halos.
The gravothermal phase $\tau = t_\mathrm{age}/t_c$ is computed using the collapse
timescale $t_c = (150/0.75)\,(\sigma_\mathrm{eff}\,\rho_s\,r_s)^{-1}(4\pi G\rho_s)^{-1/2}$
calibrated from N-body simulations \citep{Yang2023}, where $t_\mathrm{age}$ is the halo
age estimated from a quadratic fit to the Giocoli et al.\ (2011) formation redshift
relation.
The evolved SIDM density profile is evaluated using the parametric fits of Y23
(their eq.~2.3) to $\rho_s(\tau)$, $r_s(\tau)$, and $r_c(\tau)$, and the enclosed
mass within $r_{1/2}$ is computed via the analytic approximation
$M_\mathrm{SIDM}(<r_{1/2}) = \tanh(r_{1/2}/r_c(\tau)) \times M_\mathrm{NFW}(<r_{1/2};\,
\rho_s(\tau), r_s(\tau))$, which recovers the NFW result at $r_{1/2} \gg r_c$
and suppresses the enclosed mass when $r_{1/2}$ lies within the dark matter core.
We consider two SIDM benchmark models: a moderate model
($\sigma_0 = 147~\mathrm{cm^2\,g^{-1}}$, $w = 70~\mathrm{km\,s^{-1}}$) that produces
cores in the lower-mass halos, and a stronger model
($\sigma_0 = 500~\mathrm{cm^2\,g^{-1}}$, $w = 50~\mathrm{km\,s^{-1}}$) that creates
a diverse mix of cored and core-collapsed halos across the dwarf mass range.

\paragraph{Observational comparison and Rubin forecast.}
We compile a volume-complete sample of 87 Local Group dSphs within 800~kpc with Wolf
mass estimates, of which 62 have fractional mass uncertainties $\Delta M/M < 1$ and are
used for the comparison.
Enclosed mass uncertainties are dominated by statistical errors on the stellar velocity
dispersion, with median $\Delta M/M \approx 0.31$ for the SpecS5 forecast and
$\approx 0.55$ for the current literature.
We construct a differential histogram $dN/d\log_{10}M(<r_{1/2})$ in 9 equal-width
log-mass bins spanning $10^{5.5}$--$10^{8.8}~M_\odot$, with Gehrels (1986)
Poisson confidence intervals augmented by bootstrap mass measurement uncertainties
added in quadrature.
All model curves are normalised so that their total integral matches the observed
sample size, enabling a shape-only comparison; absolute normalisation requires a
full galaxy--halo connection model accounting for the luminosity function and survey
completeness, which we do not attempt here.
The current data prefer a flatter distribution than CDM at intermediate masses
($10^6$--$10^7~M_\odot$), with CDM having $\chi^2/\mathrm{dof} \approx 1.3$
compared to $\chi^2/\mathrm{dof} \approx 0.3$--0.6 for the SIDM and WDM models,
although none can be statistically rejected at $> 2\sigma$ with the current sample.

For the Rubin Observatory forecast, we estimate that $N_\mathrm{Rubin} \approx 170$
(range 150--200) dSphs within 800~kpc will have photometric detections:
$\sim$89 MW satellites detectable by Rubin within 300~kpc (Tsiane et al.\ 2025),
of which $\sim$54 are new discoveries, plus $\sim$30--50 newly discovered M31
ultra-faint satellites enabled by Rubin's southern sky coverage of the M31 system
(Mau et al.\ 2021), and $\sim$10--20 isolated Local Group dwarfs at 300--800~kpc.
We compare two scenarios for the kinematic follow-up of this photometric sample:
(i) \textit{Rubin only}, in which mass measurements are available for the 58
currently-known reliable galaxies with literature velocity dispersions, giving
Poisson errors per bin of $\sqrt{N_\mathrm{lit}/N_\mathrm{bin}} \approx 40\%$;
and (ii) \textit{Rubin + SpecS5}, in which all $\sim$170 Rubin discoveries have
SpecS5-quality velocity dispersions, reducing per-bin Poisson errors by
$\sqrt{N_\mathrm{Rubin}/N_\mathrm{lit}} \approx 1.7\times$ to $\approx 25\%$.
At this level of statistical precision, with 9 bins and model differences of $\sim$1--$2\sigma$
per bin, we forecast that the combined Rubin+SpecS5 dataset would achieve $\sim$2$\sigma$
discrimination between CDM and the moderate SIDM model integrated over the full histogram,
and $\gtrsim 3\sigma$ separation from WDM 7~keV and the primordial power-spectrum bump
model considered here, which produce qualitatively distinct features (a sharp low-mass
cutoff and an intermediate-mass excess, respectively) that are clearly distinguishable
from both CDM and SIDM even with current sample sizes.
